% ==============================================================================
% reports/sections/discussion/07_02_bias_analysis.tex
% ==============================================================================

\section{Bản chất của phép chia ngẫu nhiên và độ lệch đánh giá}

Kết quả trên Random Split với $\FOne > 99.9\%$ tương phản rõ với kết quả Time-based Split. Tuy nhiên, hai giao thức khác nhau không chỉ ở cách chia mà còn ở thành phần subtype trong Train và Test, nên phép đối chiếu này không cô lập tác động riêng của chiến lược phân tách.

\subsection{Phụ thuộc giữa các mẫu và giả thiết i.i.d.}
Hầu hết các lý thuyết nền tảng của học máy thống kê đều dựa trên giả định rằng các mẫu dữ liệu trong tập huấn luyện và kiểm thử được lấy mẫu độc lập và có cùng phân phối ($i.i.d.$):
\begin{equation}
    (\mathbf{x}_i, y_i) \stackrel{\text{i.i.d.}}{\sim} \mathcal{D}
\end{equation}

Tuy nhiên, trong an ninh mạng, các flow phát sinh từ cùng một chiến dịch tấn công có thể phụ thuộc lẫn nhau; mức độ phụ thuộc cụ thể cần được kiểm tra trên dữ liệu:
\begin{itemize}
    \item Một chiến dịch tấn công brute force kéo dài liên tục hàng giờ sử dụng một công cụ duy nhất (\texttt{patator}) chạy trên một máy trạm duy nhất.
    \item Công cụ này có cấu hình cố định về: kích thước buffer socket, thuật toán tạo luồng song song (threading model), thời gian trễ giữa các yêu cầu (delay interval) và chuỗi lệnh thử mật khẩu.
    \item Máy chủ phản hồi cũng có cấu hình cố định về các cờ giao thức và kích thước cửa sổ TCP (\texttt{Init\_Win\_bytes}).
\end{itemize}

Các đặc điểm chung của công cụ, máy trạm và cấu hình chiến dịch có thể tạo ra sự tương đồng giữa các flow cùng chiến dịch. Phân tích hiện tại không định lượng trực tiếp mức tự tương quan hoặc xác nhận rằng các flow có dấu vân tay gần như đồng nhất; đây là cơ chế hợp lý cần được kiểm chứng thêm.

\subsection{Diễn giải thận trọng về kết quả Random Split}
Random Split đưa các biến thể vào cả quá trình fit và đánh giá: Train có 1,215 mẫu SSH; Test có 1,938 mẫu SSH và 2,610 mẫu FTP. Ngược lại, Time-based Train không có SSH và Time-based Test chỉ có SSH. Do đó, hiệu năng của hai giao thức chịu ảnh hưởng đồng thời bởi khác biệt chiến lược phân tách, việc Train của Random Split đã có SSH, và phân bố subtype khác nhau trong Test.

Các flow phụ thuộc hoặc thuộc cùng chiến dịch xuất hiện ở nhiều tập có thể góp phần làm kết quả Random Split lạc quan, nhưng dữ liệu hiện có không chứng minh trực tiếp rò rỉ nhãn/đặc trưng hoặc việc mô hình ghi nhớ chiến dịch. F1 cao trong giao thức ngẫu nhiên cần được hiểu trong phạm vi thành phần subtype và phân bố Test cụ thể; để quy kết riêng tác động của cách chia, cần một thiết kế giữ tương đương việc tiếp xúc subtype và phân bố đánh giá, hoặc phân tách theo chiến dịch/khối thời gian.
